\IfFileExists{stacks-project.cls}{%
\documentclass{stacks-project}
}{%
\documentclass{amsart}
}

% For dealing with references we use the comment environment
\usepackage{verbatim}
\newenvironment{reference}{\comment}{\endcomment}
%\newenvironment{reference}{}{}
\newenvironment{slogan}{\comment}{\endcomment}
\newenvironment{history}{\comment}{\endcomment}

% For commutative diagrams we use Xy-pic
\usepackage[all]{xy}

% We use 2cell for 2-commutative diagrams.
\xyoption{2cell}
\UseAllTwocells

% We use multicol for the list of chapters between chapters
\usepackage{multicol}

% This is generally recommended for better output
\usepackage{lmodern}
\usepackage[T1]{fontenc}

% For cross-file-references
\usepackage{xr-hyper}

% Package for hypertext links:
\usepackage{hyperref}

% For any local file, say "hello.tex" you want to link to please
% use \externaldocument[hello-]{hello}
\externaldocument[introduction-]{introduction}
\externaldocument[conventions-]{conventions}
\externaldocument[sets-]{sets}
\externaldocument[categories-]{categories}
\externaldocument[topology-]{topology}
\externaldocument[sheaves-]{sheaves}
\externaldocument[sites-]{sites}
\externaldocument[stacks-]{stacks}
\externaldocument[fields-]{fields}
\externaldocument[algebra-]{algebra}
\externaldocument[brauer-]{brauer}
\externaldocument[homology-]{homology}
\externaldocument[derived-]{derived}
\externaldocument[simplicial-]{simplicial}
\externaldocument[more-algebra-]{more-algebra}
\externaldocument[smoothing-]{smoothing}
\externaldocument[modules-]{modules}
\externaldocument[sites-modules-]{sites-modules}
\externaldocument[injectives-]{injectives}
\externaldocument[cohomology-]{cohomology}
\externaldocument[sites-cohomology-]{sites-cohomology}
\externaldocument[dga-]{dga}
\externaldocument[dpa-]{dpa}
\externaldocument[sdga-]{sdga}
\externaldocument[hypercovering-]{hypercovering}
\externaldocument[schemes-]{schemes}
\externaldocument[constructions-]{constructions}
\externaldocument[properties-]{properties}
\externaldocument[morphisms-]{morphisms}
\externaldocument[coherent-]{coherent}
\externaldocument[divisors-]{divisors}
\externaldocument[limits-]{limits}
\externaldocument[varieties-]{varieties}
\externaldocument[topologies-]{topologies}
\externaldocument[descent-]{descent}
\externaldocument[perfect-]{perfect}
\externaldocument[more-morphisms-]{more-morphisms}
\externaldocument[flat-]{flat}
\externaldocument[groupoids-]{groupoids}
\externaldocument[more-groupoids-]{more-groupoids}
\externaldocument[etale-]{etale}
\externaldocument[chow-]{chow}
\externaldocument[intersection-]{intersection}
\externaldocument[pic-]{pic}
\externaldocument[weil-]{weil}
\externaldocument[adequate-]{adequate}
\externaldocument[dualizing-]{dualizing}
\externaldocument[duality-]{duality}
\externaldocument[discriminant-]{discriminant}
\externaldocument[derham-]{derham}
\externaldocument[local-cohomology-]{local-cohomology}
\externaldocument[algebraization-]{algebraization}
\externaldocument[curves-]{curves}
\externaldocument[resolve-]{resolve}
\externaldocument[models-]{models}
\externaldocument[functors-]{functors}
\externaldocument[equiv-]{equiv}
\externaldocument[pione-]{pione}
\externaldocument[etale-cohomology-]{etale-cohomology}
\externaldocument[proetale-]{proetale}
\externaldocument[relative-cycles-]{relative-cycles}
\externaldocument[more-etale-]{more-etale}
\externaldocument[trace-]{trace}
\externaldocument[crystalline-]{crystalline}
\externaldocument[spaces-]{spaces}
\externaldocument[spaces-properties-]{spaces-properties}
\externaldocument[spaces-morphisms-]{spaces-morphisms}
\externaldocument[decent-spaces-]{decent-spaces}
\externaldocument[spaces-cohomology-]{spaces-cohomology}
\externaldocument[spaces-limits-]{spaces-limits}
\externaldocument[spaces-divisors-]{spaces-divisors}
\externaldocument[spaces-over-fields-]{spaces-over-fields}
\externaldocument[spaces-topologies-]{spaces-topologies}
\externaldocument[spaces-descent-]{spaces-descent}
\externaldocument[spaces-perfect-]{spaces-perfect}
\externaldocument[spaces-more-morphisms-]{spaces-more-morphisms}
\externaldocument[spaces-flat-]{spaces-flat}
\externaldocument[spaces-groupoids-]{spaces-groupoids}
\externaldocument[spaces-more-groupoids-]{spaces-more-groupoids}
\externaldocument[bootstrap-]{bootstrap}
\externaldocument[spaces-pushouts-]{spaces-pushouts}
\externaldocument[spaces-chow-]{spaces-chow}
\externaldocument[groupoids-quotients-]{groupoids-quotients}
\externaldocument[spaces-more-cohomology-]{spaces-more-cohomology}
\externaldocument[spaces-simplicial-]{spaces-simplicial}
\externaldocument[spaces-duality-]{spaces-duality}
\externaldocument[formal-spaces-]{formal-spaces}
\externaldocument[restricted-]{restricted}
\externaldocument[spaces-resolve-]{spaces-resolve}
\externaldocument[formal-defos-]{formal-defos}
\externaldocument[defos-]{defos}
\externaldocument[cotangent-]{cotangent}
\externaldocument[examples-defos-]{examples-defos}
\externaldocument[algebraic-]{algebraic}
\externaldocument[examples-stacks-]{examples-stacks}
\externaldocument[stacks-sheaves-]{stacks-sheaves}
\externaldocument[criteria-]{criteria}
\externaldocument[artin-]{artin}
\externaldocument[quot-]{quot}
\externaldocument[stacks-properties-]{stacks-properties}
\externaldocument[stacks-morphisms-]{stacks-morphisms}
\externaldocument[stacks-limits-]{stacks-limits}
\externaldocument[stacks-cohomology-]{stacks-cohomology}
\externaldocument[stacks-perfect-]{stacks-perfect}
\externaldocument[stacks-introduction-]{stacks-introduction}
\externaldocument[stacks-more-morphisms-]{stacks-more-morphisms}
\externaldocument[stacks-geometry-]{stacks-geometry}
\externaldocument[moduli-]{moduli}
\externaldocument[moduli-curves-]{moduli-curves}
\externaldocument[examples-]{examples}
\externaldocument[exercises-]{exercises}
\externaldocument[guide-]{guide}
\externaldocument[desirables-]{desirables}
\externaldocument[coding-]{coding}
\externaldocument[obsolete-]{obsolete}
\externaldocument[fdl-]{fdl}
\externaldocument[index-]{index}

% Theorem environments.
%
\theoremstyle{plain}
\newtheorem{theorem}[subsection]{Theorem}
\newtheorem{proposition}[subsection]{Proposition}
\newtheorem{lemma}[subsection]{Lemma}

\theoremstyle{definition}
\newtheorem{definition}[subsection]{Definition}
\newtheorem{example}[subsection]{Example}
\newtheorem{exercise}[subsection]{Exercise}
\newtheorem{situation}[subsection]{Situation}

\theoremstyle{remark}
\newtheorem{remark}[subsection]{Remark}
\newtheorem{remarks}[subsection]{Remarks}

\numberwithin{equation}{subsection}

% Macros
%
\def\lim{\mathop{\mathrm{lim}}\nolimits}
\def\colim{\mathop{\mathrm{colim}}\nolimits}
\def\Spec{\mathop{\mathrm{Spec}}}
\def\Hom{\mathop{\mathrm{Hom}}\nolimits}
\def\Ext{\mathop{\mathrm{Ext}}\nolimits}
\def\SheafHom{\mathop{\mathcal{H}\!\mathit{om}}\nolimits}
\def\SheafExt{\mathop{\mathcal{E}\!\mathit{xt}}\nolimits}
\def\Sch{\mathit{Sch}}
\def\Mor{\mathop{\mathrm{Mor}}\nolimits}
\def\Ob{\mathop{\mathrm{Ob}}\nolimits}
\def\Sh{\mathop{\mathit{Sh}}\nolimits}
\def\NL{\mathop{N\!L}\nolimits}
\def\CH{\mathop{\mathrm{CH}}\nolimits}
\def\proetale{{pro\text{-}\acute{e}tale}}
\def\etale{{\acute{e}tale}}
\def\QCoh{\mathit{QCoh}}
\def\Ker{\mathop{\mathrm{Ker}}}
\def\Im{\mathop{\mathrm{Im}}}
\def\Coker{\mathop{\mathrm{Coker}}}
\def\Coim{\mathop{\mathrm{Coim}}}

% Boxtimes
%
\DeclareMathSymbol{\boxtimes}{\mathbin}{AMSa}{"02}

%
% Macros for moduli stacks/spaces
%
\def\QCohstack{\mathcal{QC}\!\mathit{oh}}
\def\Cohstack{\mathcal{C}\!\mathit{oh}}
\def\Spacesstack{\mathcal{S}\!\mathit{paces}}
\def\Quotfunctor{\mathrm{Quot}}
\def\Hilbfunctor{\mathrm{Hilb}}
\def\Curvesstack{\mathcal{C}\!\mathit{urves}}
\def\Polarizedstack{\mathcal{P}\!\mathit{olarized}}
\def\Complexesstack{\mathcal{C}\!\mathit{omplexes}}
% \Pic is the operator that assigns to X its picard group, usage \Pic(X)
% \Picardstack_{X/B} denotes the Picard stack of X over B
% \Picardfunctor_{X/B} denotes the Picard functor of X over B
\def\Pic{\mathop{\mathrm{Pic}}\nolimits}
\def\Picardstack{\mathcal{P}\!\mathit{ic}}
\def\Picardfunctor{\mathrm{Pic}}
\def\Deformationcategory{\mathcal{D}\!\mathit{ef}}
\usepackage{cleveref}
\crefname{originallemma}{original lemma}{original lemmas}
\Crefname{originallemma}{Original lemma}{Original lemmas}
\makeatletter
\newlabel{lemma-max-slice}{{04MZ}{source}{}{sourcepointer-1}{}}
\newlabel{lemma-max-slice@cref}{{[originallemma][0][]04MZ}{[1][0][]source}{}{}{}}
\makeatother
\crefname{originaldefinition}{original definition}{original definitions}
\Crefname{originaldefinition}{Original definition}{Original definitions}
\makeatletter
\newlabel{morphisms-definition-finite-type-point}{{02J1}{source}{}{sourcepointer-2}{}}
\newlabel{morphisms-definition-finite-type-point@cref}{{[originaldefinition][0][]02J1}{[1][0][]source}{}{}{}}
\makeatother
\crefname{originalreference}{original reference}{original references}
\Crefname{originalreference}{Original reference}{Original references}
\makeatletter
\newlabel{situation-slice}{{04MY}{source}{}{sourcepointer-3}{}}
\newlabel{situation-slice@cref}{{[originalreference][0][]04MY}{[1][0][]source}{}{}{}}
\makeatother
\crefname{originalremark}{original remark}{original remarks}
\Crefname{originalremark}{Original remark}{Original remarks}
\makeatletter
\newlabel{remark-local-source-apply}{{047W}{source}{}{sourcepointer-4}{}}
\newlabel{remark-local-source-apply@cref}{{[originalremark][0][]047W}{[1][0][]source}{}{}{}}
\makeatother
\crefname{originallemma}{original lemma}{original lemmas}
\Crefname{originallemma}{Original lemma}{Original lemmas}
\makeatletter
\newlabel{morphisms-lemma-locally-quasi-finite-rel-dimension-0}{{0397}{source}{}{sourcepointer-5}{}}
\newlabel{morphisms-lemma-locally-quasi-finite-rel-dimension-0@cref}{{[originallemma][0][]0397}{[1][0][]source}{}{}{}}
\makeatother
\crefname{originallemma}{original lemma}{original lemmas}
\Crefname{originallemma}{Original lemma}{Original lemmas}
\makeatletter
\newlabel{lemma-restrict-property}{{04MV}{source}{}{sourcepointer-6}{}}
\newlabel{lemma-restrict-property@cref}{{[originallemma][0][]04MV}{[1][0][]source}{}{}{}}
\makeatother
\title{Maximal slicing with locally quasi-finite stabilizer: the slice can be shrunk so that h, s' and t' are flat, finitely presented and locally quasi-finite}
\author{The Stacks Project Authors}
\date{stacks-project:more-groupoids.tex:local-snapshot-20261012; GNU FDL 1.2 or later}
\begin{document}
\maketitle
\setcounter{subsection}{1}
\noindent\textit{Editorial note.} All delivered mathematical text is the author's own, copied from the pinned Stacks Project chapter file more-groupoids.tex; no statement or proof step was rewritten, reordered or removed. The counted span is the author environment and its own proof environment. References inside the span are handed to the assembler as navigation stubs that carry the canonical tag of the cited tag only; whether each of them is load-bearing is recorded on the stub itself.\par\medskip
\noindent\textit{Allowed standard prerequisites (editorial).}\begin{itemize}
\item The maximal slicing lemma for Cohen-Macaulay groupoids at a finite type point (more-groupoids chapter, Lemma lemma-max-slice, canonical tag 04MZ; collected on this lane as L22512). cited at proof lines 2353-2354. Provides the initial slice g : U' to U with (1), (2), (3), h Cohen-Macaulay and locally of finite presentation, and dim\_\{e(u)\}(F'\_u) = dim(G'\_u).
\item Definition of a finite type point of a scheme (morphisms chapter, Definition morphisms-definition-finite-type-point, canonical tag 02J1; kind definition). cited at statement lines 2335-2336. Fixes the class of points at which the slicing problem is posed.
\item The slicing situation notation (more-groupoids chapter, Situation situation-slice, canonical tag 04MY; kind reference). cited at statement lines 2338-2339. Supplies F'\_u, G'\_u, R', s', t', h and the numbering of the conclusions.
\item Remark producing the largest open where a morphism has relative dimension at most 0 / is Cohen-Macaulay (more-groupoids chapter, Remark remark-local-source-apply, canonical tag 047W). cited at proof lines 2356-2361. Shrinks U' to the largest open U'' containing u on which h has relative dimension at most 0.
\item A morphism locally of finite presentation of relative dimension at most 0 is locally quasi-finite (morphisms chapter, Lemma morphisms-lemma-locally-quasi-finite-rel-dimension-0, canonical tag 0397). cited at proof lines 2362-2363. Turns relative dimension at most 0 into local quasi-finiteness of h.
\item Restriction lemma transferring properties of h to s' and t' of the restricted groupoid (more-groupoids chapter, Lemma lemma-restrict-property, canonical tag 04MV). cited at proof lines 2366-2368. Transfers flatness, local finite presentation and local quasi-finiteness from h to s' and t', giving (5).
\end{itemize}\medskip
\ifdefined\hypertarget\hypertarget{sourcepointer-1}{}\fi
\noindent\textit{Source reference.} Original lemma 04MZ: sources/stacks-chow/tags.txt, source lines 6060--6060. Cited at proof lines 2353-2354: provides the initial maximal slice g with its six properties.\par\medskip
\ifdefined\hypertarget\hypertarget{sourcepointer-2}{}\fi
\noindent\textit{Source reference.} Original definition 02J1: sources/stacks-chow/tags.txt, source lines 3320--3320. Cited at statement lines 2335-2336: u is a finite type point of U.\par\medskip
\ifdefined\hypertarget\hypertarget{sourcepointer-3}{}\fi
\noindent\textit{Source reference.} Original reference 04MY: sources/stacks-chow/tags.txt, source lines 6059--6059. Cited at statement lines 2338-2339: notation F'\_u, G'\_u, R', s', t', h and the numbering of the conclusions.\par\medskip
\ifdefined\hypertarget\hypertarget{sourcepointer-4}{}\fi
\noindent\textit{Source reference.} Original remark 047W: sources/stacks-chow/tags.txt, source lines 5523--5523. Cited at proof lines 2356-2361: the maximal open U'' on which h has relative dimension at most 0.\par\medskip
\ifdefined\hypertarget\hypertarget{sourcepointer-5}{}\fi
\noindent\textit{Source reference.} Original lemma 0397: sources/stacks-chow/tags.txt, source lines 4262--4262. Cited at proof lines 2362-2363: relative dimension at most 0 plus local finite presentation gives local quasi-finiteness.\par\medskip
\ifdefined\hypertarget\hypertarget{sourcepointer-6}{}\fi
\noindent\textit{Source reference.} Original lemma 04MV: sources/stacks-chow/tags.txt, source lines 6056--6056. Cited at proof lines 2366-2368: transfers flatness, local finite presentation and local quasi-finiteness from h to s' and t'.\par\medskip
\begin{lemma}
\label{lemma-max-slice-quasi-finite}
Let $S$ be a scheme.
Let $(U, R, s, t, c, e, i)$ be a groupoid scheme over $S$.
Let $G \to U$ be the stabilizer group scheme.
Assume $s$ and $t$ are Cohen-Macaulay and locally of finite presentation.
Let $u \in U$ be a finite type point of the scheme $U$, see
Morphisms, Definition \ref{morphisms-definition-finite-type-point}.
Assume that $G \to U$ is locally quasi-finite.
With notation as in
Situation \ref{situation-slice}
there exist an affine scheme $U'$ and a morphism $g : U' \to U$ such that
\begin{enumerate}
\item $g$ is an immersion,
\item $u \in U'$,
\item $g$ is locally of finite presentation,
\item the morphism $h : U' \times_{g, U, t} R \longrightarrow U$
is flat, locally of finite presentation, and locally quasi-finite, and
\item the morphisms $s', t' : R' \to U'$ are flat,
locally of finite presentation, and locally quasi-finite.
\end{enumerate}
\end{lemma}

\begin{proof}
Take $g : U' \to U$ as in
Lemma \ref{lemma-max-slice}.
Since $h^{-1}(u) = F'_u$ we see that $h$ has relative dimension
$\leq 0$ at $(u, e(u))$. Hence, by
Remark \ref{remark-local-source-apply},
we obtain an open subscheme $U'' \subset U'$ such that
$u \in U''$ and $U'' \times_{g, U, t} R$ is the maximal open subscheme
of $U' \times_{g, U, t} R$ on which $h$ has relative dimension $\leq 0$.
After replacing $U'$ by $U''$ we see that $h$ has relative dimension $\leq 0$.
This implies that $h$ is locally quasi-finite by
Morphisms, Lemma \ref{morphisms-lemma-locally-quasi-finite-rel-dimension-0}.
Since it is still locally of finite presentation and Cohen-Macaulay we see
that it is flat, locally of finite presentation and locally quasi-finite,
i.e., (4) above holds. This implies that $s'$ is flat, locally of finite
presentation and locally quasi-finite as a base change of $h$, see
Lemma \ref{lemma-restrict-property}.
\end{proof}

\end{document}
